A belief $\belief_t$ is a hypothesis about the environment that the agent carries between decisions. It yields a prediction, evidence $o_{t+1}$ returns, a \emph{test relation} supports or falsifies the hypothesis, and a revision produces $\belief_{t+1}$ for the next action (Fig.~\ref{fig:loop}). The test may be a symbolic verifier, a consistency check, or a comparison with an external record. Partially observed decision-making gives this cycle its most principled realisation \citep{astrom1965optimal,kaelbling1998planning}. With hidden state $z_t$ and history $\hist_t$, the belief is $\belief_t = p(z_t \mid \hist_t)$, and a Bayes filter carries it forward as
\begin{equation}
\belief_{t+1}(z') \propto p(o_{t+1} \mid z') \textstyle\sum_{z} p(z' \mid z, a_t)\, \belief_t(z)
\eqfilterlab
\end{equation}
\citep{thrun2005probabilistic}. Its three factors are the Bayesian form of the three parts the review follows. The factor $\belief_t$ is the \emph{belief state}, what is believed and how firmly, the transition $p(z' \mid z, a_t)$ is the \emph{revision} that carries it across the action, and the likelihood $p(o_{t+1} \mid z')$ is the \emph{test}. Language-model agents rarely separate the filter's prediction and correction steps, so we call revision the whole move from $\belief_t$ to $\belief_{t+1}$, which the test feeds. A verifier or consistency check is a non-probabilistic test of the same kind and is coded as one. Agentic reinforcement learning casts the language-model agent as such a POMDP, with the context as history and any explicit state in a working memory the model writes \citep{zhang2025landscape,sumers2023cognitive,yao2023react}, and we write $\belief_t$ for whatever explicit state the span carries before the action $a_t$.

We call such an object a \emph{belief state} only if (i) it is a claim about the environment carried between decisions, (ii) it carries, or admits, uncertainty, and (iii) it is \emph{testable}, committing to a claim that an observation the agent receives can support or falsify, usually a prediction of the next one. Conditions (i) and (ii) are properties of the belief state and fix the operation \op{Construct}. Condition (iii) requires a test and fixes \op{Test}, which makes the state a hypothesis rather than a record. \op{Revise} is the revision, which belief-revision theory constrains to change the claim without discarding the rest \citep{alchourron1985logic,friston2010free}. We therefore place every system by which of the three parts it has and where each comes from, whether written by the model, fixed by the designer, learned from data, or absent. A trace \emph{memory} fails (i) and (ii). A store that attaches a probability and an update rule to each item meets them but scores nothing against what returns, so it fails (iii) \citep{luo2026storage}. A \emph{world model} is itself a revisable hypothesis, a learned revision and test whose state the agent cannot read, so it meets (iii) but seldom (i) \citep{li2026bridging}.
